% Moduł wspólny, wykorzystywany w geometrii i topologii algebraicznej.
% Dwa historyczne identyfikatory odsyłaczy odpowiadają jednemu twierdzeniu.
% Guard zapobiega podwójnemu składowi w książce łączącej oba rozdziały.
\ifcsname biblioteka@slabelefschetz@loaded\endcsname
Twierdzenie wraz z dowodem zostało już sformułowane wcześniej
(w twierdzeniu~\ref{tw:ga10-weak}).
\else
\expandafter\gdef\csname biblioteka@slabelefschetz@loaded\endcsname{}
\begin{twierdzenie}[Słabe twierdzenie Lefschetza]
\label{tw:ga10-weak}\label{thm:xvii-weak-lefschetz}
Niech $X\hookrightarrow\mathbb P^N_{\mathbb C}$ będzie gładką,
spójną odmianą projektową nad $\mathbb C$ o wymiarze zespolonym
$n\ge 2$. Niech $H\subset\mathbb P^N_{\mathbb C}$ będzie hiperpłaszczyzną,
a $Y=X\cap H$ niepustym gładkim przekrojem, z zanurzeniem
$i\colon Y\hookrightarrow X$. Wówczas homomorfizm kohomologii singularnej
\[
 i^*\colon H^q(X^{\mathrm{an}};\mathbb Z)
 \longrightarrow H^q(Y^{\mathrm{an}};\mathbb Z)
\]
jest izomorfizmem dla $q<n-1$ oraz monomorfizmem dla $q=n-1$.
Wersja kohomologiczna z dowolnymi współczynnikami w ciele
wynika z tej samej konstrukcji. W szczególności $Y$ jest spójna. Dla odpowiednio ustalonych
punktów bazowych indukowane odwzorowanie
$\pi_j(Y^{\mathrm{an}})\to\pi_j(X^{\mathrm{an}})$
jest izomorfizmem dla $1\le j<n-1$ oraz epimorfizmem dla $j=n-1$.
\end{twierdzenie}
\begin{proof}[Redukcja do geometrycznego twierdzenia Lefschetza]
Korzystamy z następującego \emph{głębokiego wyniku zewnętrznego}:
przestrzeń $X^{\mathrm{an}}$ ma typ homotopijny przestrzeni
otrzymanej z $Y^{\mathrm{an}}$ przez dołączenie komórek
rzeczywistego wymiaru co najmniej $n$. Dowód tego wyniku
wymaga zespolonej teorii Morse'a i funkcji plurisubharmonicznych
(lub odpowiedniej teorii znikania). Nie wyprowadzamy go z samej
gładkości $X$.

Względny kompleks komórkowy pary $(X^{\mathrm{an}},Y^{\mathrm{an}})$
nie ma komórek w stopniach mniejszych niż $n$, więc
$H_q(X^{\mathrm{an}},Y^{\mathrm{an}};\mathbb Z)=0$ dla $q<n$.
Względna aproksymacja komórkowa daje również
$\pi_j(X^{\mathrm{an}},Y^{\mathrm{an}})=0$ dla $j<n$.
Długi ciąg dokładny homotopii pary zapewnia stąd podane
izomorfizmy i epimorfizm. Należy przy tym pamiętać, że
w niskich stopniach grupy homotopii wymagają właściwej
interpretacji z ustalonymi punktami bazowymi.

Analogicznie $H^q(X^{\mathrm{an}},Y^{\mathrm{an}};\mathbb Z)=0$
dla $q<n$. W ciągu dokładnym kohomologii pary
\[
 \cdots\longrightarrow H^q(X,Y;\mathbb Z)
 \longrightarrow H^q(X;\mathbb Z)
 \xrightarrow{i^*} H^q(Y;\mathbb Z)
 \longrightarrow H^{q+1}(X,Y;\mathbb Z)\longrightarrow\cdots
\]
(przez $X,Y$ w tym ciągu rozumiemy przestrzenie analityczne)
dla $q<n-1$ znikają obie grupy względne sąsiadujące z $i^*$,
a dla $q=n-1$ znika grupa z lewej. Otrzymujemy odpowiednio
izomorfizm i monomorfizm. Ponieważ $n\ge2$, dla $q=0$
otrzymujemy izomorfizm $H^0(X;\mathbb Z)\simeq H^0(Y;\mathbb Z)$,
a spójność $X$ daje spójność $Y$.
\end{proof}
\fi
